\metatitle{Unit interval cographs and OEIS A152947}
\metadescription{Classification and enumeration of natural unit interval graphs which are cographs.}


\section[oeisProofA152947]{Unit interval cographs and A152947}

This page is part of the \hyperref[oeisProofs]{OEIS proof collection}.

A \defin{cograph} is a graph with no induced path on four vertices.
We count cographs among the natural
\hyperref[chromaticQuasisymmetricUnitIntervalGraph]{unit interval graphs}
encoded by area sequences.

\begin{theorem}
Let $c_n$ be the number of connected area sequences of length $n$ whose
unit interval graph is a cograph.  For $n \geq 1$,
\[
c_n=1+\binom{n-1}{2},
\]
so $(c_n)_{n\geq1}$ is \oeis{A152947}.

More precisely, a connected unit interval graph is a cograph if and only if
it is a complete graph or has the form
\[
K_r \mathbin{\vee} (K_p \mathbin{\sqcup} K_q),
\qquad p,q,r\geq1,
\]
where $\vee$ denotes the graph join.
\end{theorem}

\begin{proof}
We use the recursive characterization of cographs: every connected cograph
with at least two vertices is a nontrivial join of smaller cographs.  Refine
such a decomposition into join-indecomposable factors.  If two factors both
contain a nonedge, then a nonadjacent pair from each factor induces a
$4$-cycle, since all edges between different factors are present.  This is
impossible because unit interval graphs are chordal.  Thus at most one factor
is not complete.  The complete factors are single vertices in the refined
decomposition; together they form a nonempty clique $K_r$ of universal
vertices.

Let $H$ be the remaining noncomplete factor, if it exists.  It is a
join-indecomposable cograph, so it is disconnected.  Unit interval graphs are
claw-free.  Hence $H$ has independence number at most two: otherwise a
universal vertex together with three independent vertices of $H$ would induce
a claw.  A maximum independent set of a disconnected graph contains a vertex
from every component.  It follows that $H$ has exactly two components and
that each component is complete.  Consequently
$H=K_p\sqcup K_q$ for some $p,q\geq1$.  This proves the claimed structural
form.  Conversely, complete graphs are cographs, and cographs are closed
under disjoint union and join.  The displayed graphs are also unit interval
graphs, as the following area sequences show.

In a unit interval representation of
$K_r\vee(K_p\sqcup K_q)$, the two nonadjacent clique classes occur on
opposite sides of the universal class.  Choosing which class is on the left,
the area sequence is
\[
(0,1,\dotsc,p+r-1,r,r+1,\dotsc,r+q-1).
\]
The complete graph $K_n$ has area sequence $(0,1,\dotsc,n-1)$.  Conversely,
these sequences produce exactly the displayed graphs.  We therefore obtain
one complete case and one case for every ordered composition
$n=p+r+q$ into three positive parts.  There are
$\binom{n-1}{2}$ such compositions, proving the formula.
\end{proof}

There is also a simple enumeration without the connectedness condition.
Every area sequence decomposes uniquely into consecutive connected blocks.
Thus, if $a_n$ counts all area sequences of length $n$ whose graph is a
cograph, then
\[
C(x)\coloneqq\sum_{n\geq1}c_nx^n
=\frac{x}{1-x}+\frac{x^3}{(1-x)^3}
\]
and
\[
\sum_{n\geq0}a_nx^n
=\frac{1}{1-C(x)}
=\frac{(1-x)^3}{1-4x+5x^2-3x^3}.
\]
The first values of $(a_n)_{n\geq0}$ are
\[
1,1,2,5,13,33,82,202,497,1224,3017,7439,18343.
\]

% Related Rust: combinatoric-core/src/graph.rs, Graph::is_p4_free.

